% frontmatter/english/como-usar.tex -- "How to use this book": las
% instrucciones para el adulto de frontmatter/como-usar.tex, en inglés.
\section*{How to use this book}

This book is for a child of five or six who is starting out with
numbers: recognising them, counting them and writing them -- and, as the
year goes on, comparing them, adding and taking away. It belongs to the
family of \emph{First Words} and \emph{Read and Draw}: the same
characters (Lucía, her little brother Dani, their dog Toby, Mum, Dad and
Grandma Rosa). And it is, page by page, the English edition of
\emph{Aprendo los números}: the same number, the same pictures and the
same activity every day, so a child can use either book, or both.

There is a page for every weekday of the year, Monday to Friday,
holidays included. Five or ten minutes a day is plenty. Every page works
the same way:

\begin{itemize}[leftmargin=6mm]
\item \textbf{Date}, \textbf{Done} and \textbf{Notes}: for the grown-up.
  Ticking each day whether it was done alone, with help or with
  difficulty takes a second, and by the end of the year it is a record
  of everything that has been learnt.
\item \textbf{Today's number}, in the green box: the big number, its
  name, the \textbf{ten frame} -- two rows of five squares, with a dot
  for each one, so you can see at a glance how much a number is (5 is a
  whole row; 7 is a row and two more) -- and that many things, drawn.
  Underneath, a sentence from the story, which the grown-up reads out:
  today's number, at Lucía's house.
\item An \textbf{activity}, which nearly always ends in colouring,
  circling or drawing. The instructions (in small grey print) are
  \textbf{read out by a grown-up}.
\end{itemize}

Each activity is explained, one by one, on the next page.

\subsection*{How to help}

{\small
\begin{itemize}[leftmargin=6mm, itemsep=1pt, topsep=2pt]
\item Counting is touching: point to each thing with a finger as you
  count it, and only once. If one gets skipped or counted twice, count
  together, slowly.
\item Say the number's name every time it is traced or found: the
  number is seen, said and written.
\item If something is hard, it can be done again the next day with
  things from around the house: spoons, buttons, shoes.
\item In English, 13 and 30, 14 and 40... sound almost the same. When
  they come up, say the end of the word clearly: thir\textbf{teen},
  \textbf{thir}ty.
\item At the end of each part there is a medal, and at the end of the
  book, a diploma and the \textbf{answers} for the pages that have one.
\end{itemize}}
\newpage
